\scanpage{78}
\mspage{\handwritten{- 1 -}}

\noindent\uline{F\,O\,R\,E\,W\,O\,R\,D}

\medskip
\noindent\begin{minipage}{\linewidth}\small
I.\qquad
% One continuous oblique pen stroke, from above the opening down into the footnote.
\tikz[remember picture,overlay]{\coordinate (cotesixtythreestart) at (-0.7em,1.8ex);}
The principal aim of the present Seminar is to develop the formalism of the ``Weil cohomology'' of Schemes. Starting essentially from the results which are proved here, well-known arguments, moreover due to Weil himself, \'{p}ermit one to deduce part of the Weil conjectures on the $\mathrm{L}$-functions of non-singular projective varieties over a finite field\textsuperscript{\handwritten{(1)}}. This deduction will moreover be presented in SGA \handwritten{1964/65}, in a seminar which will follow this one.

In the present seminar, we restrict ourselves to the study of the cohomology of schemes with respect to the \uline{étale topology} . This topology by its description is very close to the topologies of the usual topological manifolds\handwritten{,} and we shall see that most of the classical results concerning the cohomology of ordinary.topological spaces (various spectral sequences, finiteness theorems, K\handwritten{\"}{u}nneth, duality, Lefschetz theorem) can be formulated and proved in the new context, provided one restricts oneself if necessary to torsion sheaves \uline{prime to the residue characteristics} of the schemes considered\unsure{.|,} We shall obtain a cohomological theory ``with coefficient of characteristic $0$'' (as requested by Weil) by an e\makebox[0pt][l]{\deleted{ss}}ssentially trivial passage to the inverse limit (cf SGA \handwritten{1965}), allowing one to define a cohomology with coefficient \deleted{\unsure{of}}\handwritten{\unsure{i|i}n} the ring $\underline{\mathrm{Z}_{\mathord{\rlap{\handwritten{$\scriptstyle/$}}\mathrm{l}}}}$ of $\mathord{\rlap{\handwritten{$/$}}\mathrm{l}}$-adic integers starting from the coefficients $\underline{\mathrm{Z}}\handwritten{/}\mathord{\rlap{\handwritten{$/$}}\mathrm{l}}^{\handwritten{\nu}}\underline{\mathrm{Z}}$ ., $\handwritten{\nu\longrightarrow}+\handwritten{\infty}$ . When $\mathord{\rlap{\handwritten{$/$}}\mathrm{l}}$ is prime to the residue characteristics, this cohomology possesses all the good properties usual in classical cohomology with coefficient $\underline{\mathrm{Z}}$ (and therefore lends itself to the formulation of the Weil conjectures). Moreover, when the schemes considered are of finite type over the field of complex numbers $\underline{\mathrm{C}}$, we shall see that the cohomology (with torsion coefficient\handwritten{s}) is essentially identical to the cohomology

\smallskip
\noindent\rule{0.32\linewidth}{0.4pt}\par
{\footnotesize\noindent\textsuperscript{\handwritten{(1)}}\quad At the time of writing these lines, it is not proved that the eigenvalues of the Frobenius homomorphism acting on the $\mathrm{H}^{\handwritten{i}}\handwritten{(}\mathrm{X},\underline{\mathrm{Z}}_{\handwritten{\ell}}\handwritten{)}$ are algebraic integers, nor a fortiori that their\handwritten{\tikz[remember picture,overlay]{\draw[line width=0.4pt] (cotesixtythreestart) -- (-0.45em,0.15ex);}} absolute values are equal to $\mathrm{q}^{\handwritten{i/2}}$ .\par}
\end{minipage}
% archive sheet number: 77
